\documentclass[10pt,letterpaper,unboxed,cm, fleqn]{article}
\usepackage[margin=1in]{geometry}
\usepackage{graphicx}
\usepackage{enumerate}
\usepackage{enumitem}
\usepackage{amsmath}
\usepackage{amssymb}
\usepackage{float}
\usepackage{mathtools}
\usepackage{empheq}

\newcommand{\st}{~\mid~}
\newcommand{\ind}{$~~~$}
\usepackage{xcolor}

\usepackage{fancyhdr}

\fancyhf{}
\pagestyle{fancy}

\everymath{\displaystyle}

\setlength{\mathindent}{0pt}

\begin{document}

\begin{tabular}{|c|c|c|c|c|c|c|c|c|}
    \hline
    & 3/L0 & 12/L1 & 13/L2 & 14/L3 & 15/L4 & 16/L5 & 17/L6 & 18/L7 \\
    \hline
    1/H0 & DOWN & ADV & & & & RIGHT & LEFT & UP \\
    \hline
    2/H1 & CORRECT & n & h & y & 6 & BKSP & j & . \\
    \hline
    4/H2 & WORD ERASER & b & g & t & 5 & = & RET & , \\
    \hline
    5/H3 & SPC & v & f & r & 4 & - & 1/2 & m \\
    \hline
    6/H4 & MARGIN & c & d & e & 3 & o & p & ' \\
    \hline
    7/H5 & CODE & x & s & w & 2 & 9 & 0 & ; \\
    \hline
    8/H6 & / & z & a & q & 1 & 8 & i & l \\
    \hline
    9/H7 & RSHIFT & LSHIFT & CAPS & TAB & TABs & 7 & u & k \\
    \hline
\end{tabular}

\end{document}

